%% conference_review.tex — конференционная версия обзора
%% «Чувствительные показатели здоровья почвы».
%% Источник: paper/review_soil_health_indicators.md (сокращение и сопоставление).
%% Сборка: pdflatex conference_review.tex (дважды).

\documentclass{ai4science_ysda}

\InformationInRussian

\TITLE{Чувствительные показатели здоровья почвы: сравнение, условия
применения и ограничения}
\ShortTITLE{Чувствительные показатели здоровья почвы}
\AUTHOR{Анонимные авторы}

\Abstract{Здоровье почвы описывают через физические, химические и
биологические свойства, поэтому выбор индикаторов для мониторинга остаётся
предметом обсуждения. В работе сопоставлены три группы показателей:
физико-химические (органический углерод и его нормированные формы),
лабильные фракции углерода и микробиологические показатели. Для каждой
группы указаны воздействия, при которых наблюдается отклик, и ограничения
интерпретации. В рассмотренных работах лабильные фракции отзывчивы к
сокращённой обработке почвы и внесению органического вещества,
микробиологические показатели реагируют на более широкий набор
агротехнических практик, а физико-химические показатели служат основой
оценки при нормировке по педо-климатическому контексту. Приведена
сравнительная таблица с условиями применения и ограничениями.}

\KEYWORDS{здоровье почвы, индикаторы, органический углерод, лабильные
фракции углерода, микробиологические показатели}

\begin{document}

\section*{Введение и постановка задачи}

\porozhdeno{} Здоровье почвы определяют как её способность функционировать
как живая система и обеспечивать экосистемные услуги. Для мониторинга
требуются показатели, которые реагируют на изменение управления и при этом
сохраняют связь с долгосрочными изменениями состояния почвы.

Органический углерод почвы (SOC) входит в число наиболее широко
используемых показателей наряду с pH и доступными фосфором и калием
\cite{bunemann2018}. Абсолютное содержание SOC зависит от климата, типа
почвы и землепользования, что затрудняет назначение целевых значений.
Задача работы — сопоставить показатели трёх групп, указать воздействия, при
которых каждый из них даёт отклик, и ограничения, при которых вывод
перестаёт быть верным.

\section{Подготовка обзора}

\porozhdeno{} Обзор подготовлен в ходе курса «ИИ-ассистенты для
исследователя». Основная задача проекта — разработка агента для работы с
научной литературой и подготовки обзора показателей здоровья почвы. Полная
версия обзора основана на корпусе из 47 научных статей; для конференции
полный текст сокращён. В сокращённой версии сохранены группы индикаторов,
их сравнение, условия применения и ограничения. Показатели сопоставлены по
воздействию, отклику, условиям исследования и ограничениям интерпретации.

Число статей корпуса приведено по данным проекта. Повторный подсчёт статей
и повторное чтение корпуса при подготовке этой версии не выполнялись.
Формальный протокол систематического поиска не применялся, поэтому полнота
корпуса не оценена.

\section{Физико-химические показатели и нормировка SOC}

\porozhdeno{} Нормировка содержания углерода по содержанию глины основана
на связи SOC с тонкими частицами \cite{feeney2024}. Отношение SOC/clay
используется как порог насыщения углеродом, однако нормировка только по
глине смещена по текстуре и климату и может давать вводящие в заблуждение
оценки структуры почвы и потенциала накопления углерода \cite{feeney2024}.

Альтернативой служит индикатор O/T SOC — отношение наблюдаемого содержания
углерода к типичному для данной педо-климатической зоны \cite{feeney2024}.
Типичные значения получены по данным программы LUCAS (2009--2018,
19\,855 минеральных верхних горизонтов 0--20 см): выделены 12
педо-климатических зон с типичными концентрациями 5,99--39,65~г/кг
\cite{feeney2024}. Индикатор менее чувствителен к содержанию глины,
землепользованию и климату и лучше отражает различия в пористости и
запасах углерода \cite{feeney2024}. Физико-химические показатели реагируют
медленнее биологических, поэтому их применяют в сочетании с другими типами
показателей \cite{cardoso2013}.

\section{Лабильные фракции углерода}

\porozhdeno{} Изменения, вызванные агротехническими практиками, часто
трудно обнаружить по общему содержанию SOC, поэтому измеряют быстро
меняющиеся пулы \cite{bongiorno2019}. В десяти европейских долгосрочных
полевых экспериментах сравнены пять фракций: растворённый органический
углерод (DOC), его гидрофильная часть (Hy-DOC), углерод, окисляемый
перманганатом (POXC), углерод, экстрагируемый горячей водой (HWEC), и
углерод органического вещества в виде частиц (POMC) \cite{bongiorno2019}.

Сокращённая обработка почвы и высокое внесение органического вещества
повышали концентрации лабильных фракций по сравнению с традиционной
обработкой и низким внесением. К обоим факторам наиболее отзывчивыми
оказались POXC и POMC; POXC дополнительно наиболее тесно коррелировал с
химическими, физическими и биологическими показателями
\cite{bongiorno2019}. Измерение POXC недорого и выполняется быстро.

Отклик зависит от условий отбора проб: часть исследований не обнаружила
эффекта обработки и удобрения на лабильные фракции, а химический состав и
сезонная динамика POXC изучены слабо \cite{bongiorno2019}.

\section{Микробиологические показатели}

\porozhdeno{} Метаанализ 166 статей оценил отклик активных пулов
органического вещества и микробиологических показателей на шесть практик:
внесение экзогенного органического вещества (EOM), сохранение растительных
остатков (RR), покровные культуры (CC), диверсификацию севооборота (CR),
сокращённую обработку почвы (TR) и переход к инновационным или органическим
системам (IOS) \cite{cusset2024}. Климатические факторы значимо не влияли
на отклик, тогда как содержание глины и органического вещества, pH,
длительность опыта и глубина отбора проб его модулировали
\cite{cusset2024}.

К EOM, RR, CC и TR были чувствительны как физико-химические (SOC, общий
азот STN), так и микробиологические показатели. В частности, покровные
культуры повышали SOC на 7\,\%, STN на 8\,\%, микробную биомассу углерода
(MBC) на 14\,\%, общие PLFA на 19\,\% и $\beta$-глюкозидазу (B-GLU) на
11\,\% \cite{cusset2024}. Диверсификация севооборота и инновационные
системы не влияли на SOC и STN, но повышали MBC на 7--34\,\%, PLFA на
8--24\,\%, а B-GLU — на 58\,\% при переходе к инновационным системам
\cite{cusset2024}. Величины относятся к перечисленным практикам и условиям
метаанализа.

Ферментативная активность служит самостоятельным биоиндикатором:
дегидрогеназы, фосфатазы и $\beta$-глюкозидазы отражают влияние обработки,
органических удобрений и севооборота, однако методы чувствительны к
пробоподготовке и условиям хранения образцов \cite{bhupenchandra2026}.
Функциональные гены циклов углерода, азота и фосфора чувствительны к
вариантам удобрения в десятилетнем полевом опыте \cite{jia2025}, но требуют
молекулярно-биологических методов. Состав микробных сообществ реагирует на
управление, а величина эффекта зависит от биома \cite{trivedi2016}.
Интерпретацию осложняют прайминг-эффекты, при которых микробная биомасса
выступает движущей силой оборота углерода и азота \cite{kuzyakov2010}, а
выбор биоиндикаторов ограничен их временным и пространственным разрешением
\cite{bhaduri2022}.

\section{Сравнение и условия применения}

\porozhdeno{} Сопоставление рассмотренных работ позволяет сформулировать
три вывода об условиях применения. Физико-химические показатели служат
базой оценки органического статуса почвы, но откликаются медленно и требуют
нормировки по педо-климатическому контексту \cite{feeney2024,cardoso2013}.
Лабильные фракции отзывчивы к сокращённой обработке почвы и внесению
органического вещества \cite{bongiorno2019}. Микробиологические показатели
реагируют на более широкий набор практик, включая те, которые не меняют SOC
и STN в краткосрочных опытах \cite{cusset2024}.

Формулировки о наибольшей чувствительности относятся к конкретным работам и
воздействиям. В опытах Bongiorno и соавторов наиболее отзывчивыми к
обработке и удобрению оказались POXC и POMC \cite{bongiorno2019}; в
метаанализе Cusset и соавторов микробиологические показатели проявили
отклик к более широкому спектру практик \cite{cusset2024}. Единый рейтинг
показателей по этим данным не строится.

Согласованность выводов с более широкой литературой неполная. На одной
долгосрочной площадке снижение обработки не изменило содержание
органического вещества, ёмкость катионного обмена и POXC в традиционной
системе \cite{pearsons2023}. В сводке долгосрочных опытов Европы и Китая
внесение органического вещества сильнее всего влияло на численность
дождевых червей, затем на урожайность, содержание органического вещества и
агрегатную стабильность; отказ от обработки (no-tillage) повышал
агрегатную стабильность и содержание органического вещества в верхних
горизонтах при снижении урожайности \cite{bai2018}.

Для масштабного мониторинга предложен минимальный набор из трёх
показателей: SOC, агрегатная стабильность и потенциал минерализации
углерода за 24 часа \cite{bagnall2023}.

\begin{table}[!htbp]
  \centering
  \caption{Сравнение индикаторов: воздействия, отклик и ограничения
  (обозначения практик --- раздел 4).}
  \label{tab:indicators}
  \small
  \setlength{\tabcolsep}{2pt}
  \renewcommand{\arraystretch}{0.9}
  \begin{tabular}{@{}>{\raggedright\arraybackslash}p{0.17\textwidth}
      >{\raggedright\arraybackslash}p{0.21\textwidth}
      >{\raggedright\arraybackslash}p{0.18\textwidth}
      >{\raggedright\arraybackslash}p{0.19\textwidth}
      >{\raggedright\arraybackslash}p{0.13\textwidth}@{}}
    \hline
    \textbf{Индикатор} & \textbf{Воздействие или условия} &
    \textbf{Характер отклика} & \textbf{Ограничение} &
    \textbf{Источник} \\
    \hline
    SOC, STN & EOM, RR, CC, TR & рост; слабый отклик на CR & медленный;
    климат и текстура & \cite{cusset2024} \\
    SOC/clay & нормировка по глине & порог насыщения углеродом & смещение по
    текстуре и климату & \cite{feeney2024} \\
    O/T SOC & педо-климатические зоны & различает пористость и запасы
    углерода & не для органических почв и подпочв & \cite{feeney2024} \\
    POXC & TR, EOM & рост при TR и высоком EOM & зависит от условий и срока
    отбора & \cite{bongiorno2019} \\
    POMC & TR, EOM & среди наиболее отзывчивых & неоднородность остатков &
    \cite{bongiorno2019} \\
    DOC, HWEC & TR, EOM & слабее, чем у POXC и POMC & зависит от свойств
    почвы и сезона & \cite{bongiorno2019} \\
    MBC & EOM, RR, CC, CR, IOS & рост, в том числе без роста SOC & стоимость
    и квалификация & \cite{cusset2024} \\
    PLFA & EOM, RR, CC, CR, IOS & рост при CR и IOS & сложный протокол &
    \cite{cusset2024} \\
    Ферменты (B-GLU, FDA, URE) & EOM, RR, CC, CR, IOS & рост; максимум B-GLU
    при IOS & стандартизация анализов & \cite{cusset2024,bhupenchandra2026} \\
    Дождевые черви & EOM, no-tillage, травы в ротации & рост численности &
    меньше данных; пестициды & \cite{bai2018} \\
    Агрегатная стабильность & no-tillage, органическое земледелие & рост &
    зависит от глубины и текстуры & \cite{bai2018} \\
    \hline
  \end{tabular}
\end{table}

\section{Проверка утверждений}

\porozhdeno{} Сокращённый текст сопоставлен с исходным обзором по составу
разделов, набору цитируемых работ и численным эффектам. Сборка выполнена
дважды командой \texttt{pdflatex}; проверены число страниц, неопределённые
ссылки и переполнения строк. Техническая проверка проекта
(\texttt{python checks/check.py}) относится к целевым файлам основного
обзора и содержание конференционного текста не подтверждает.

Сверка конференционного текста с первоисточниками по PDF не выполнялась.
Отчёты выборочной сверки в каталоге \texttt{reports/} относятся к прежним
версиям обзора и проверкой настоящего текста не являются, поэтому число
подтверждённых утверждений не приводится.

\section{Ограничения}

\porozhdeno{} Корпус обзора включает 47 научных статей; полнота охвата
литературы по формальному протоколу систематического поиска не
оценивалась. Сопоставимость исследований ограничена различиями
почвенно-климатических условий, длительности опытов, практик управления и
методов измерения. Индикатор O/T SOC разработан для минеральных верхних
горизонтов; для органических почв и подпочвенных горизонтов набор типичных
концентраций не определён, а классы чувствительны к точности цифровых карт
\cite{feeney2024}. Объединение съёмок LUCAS исключает анализ динамики во
времени \cite{feeney2024}.

Лабильные фракции зависят от условий отбора проб и сезонной динамики
\cite{bongiorno2019}. Микробиологические показатели требуют больше времени,
квалификации и стоимости \cite{cusset2024}; в метаанализе использовалась
фиксированная глубина отбора \cite{cusset2024}. Эффекты микробных
показателей зависят от биома \cite{trivedi2016}, функциональные гены
требуют специализированной инфраструктуры \cite{jia2025}, разрешение
биоиндикаторов ограничивает их перенос между условиями \cite{bhaduri2022}.
Оценки чувствительности качественные; единая шкала не вводится.

\section*{Заключение}

\porozhdeno{} Набор показателей выбирают по воздействию, отклик на которое
нужно обнаружить, по доступным ресурсам и по педо-климатическим условиям
площадки. Показатели с медленным откликом дают оценку состояния, показатели
с быстрым откликом — сигнал об изменении практики до того, как оно
проявится в общем содержании углерода. Единый рейтинг показателей по
рассмотренным работам не строится: оценки отклика привязаны к конкретным
опытам, глубинам отбора и срокам наблюдений. Дальнейшая работа —
стандартизация протоколов измерения и проверка показателей в разных
педо-климатических зонах.

\begin{aiusage}[Пометки: \porozhdeno{} — фрагмент порождён моделью,
\istochnik{} — фрагмент взят из источника.]
  Аннотация, разделы 1--7 & \porozhdeno{} & open-weight LLM через Yandex AI
    Studio; версия не зафиксирована & сопоставлено с исходным обзором;
    редактура языка моделью; сверка по PDF не выполнялась \\
  Таблица~\ref{tab:indicators} & \porozhdeno{} & тот же инструмент & по
    таблице §7 исходного обзора \\
  Библиография & \istochnik{} & перенос \texttt{bibitem} из исходного .tex &
    метаданные не сверялись \\
  Сборка и проверка & \porozhdeno{} & pdflatex, checks/check.py & две
    сборки, код в отчёте \\
\end{aiusage}

\begin{thebibliography}{13}
\bibitem{feeney2024}
Feeney, C.J., Bentley, L., De Rosa, D., Panagos, P., Emmett, B.A., Thomas, A.,
Robinson, D.A. (2024). Benchmarking soil organic carbon (SOC) concentration
provides more robust soil health assessment than the SOC/clay ratio at European
scale. \emph{Science of the Total Environment}, 951, 175642.

\bibitem{bai2018}
Bai, Z., Caspari, T., Ruiperez Gonzalez, M., Batjes, N.H., Mäder, P.,
Bünemann, E.K., et al. (2018). Effects of agricultural management practices on
soil quality: A review of long-term experiments for Europe and China.
\emph{Agriculture, Ecosystems and Environment}, 265, 1--7.

\bibitem{bongiorno2019}
Bongiorno, G., Bünemann, E.K., Oguejiofor, C.U., Meier, J., Gort, G., Comans, R.,
Mäder, P., Brussaard, L., de Goede, R. (2019). Sensitivity of labile carbon
fractions to tillage and organic matter management and their potential as
comprehensive soil quality indicators across pedoclimatic conditions in Europe.
\emph{Ecological Indicators}, 99, 38--50.

\bibitem{cusset2024}
Cusset, E., Bennegadi-Laurent, N., Recous, S., Bernard, P.-Y., Perrin, A.-S.,
Tscheiller, R., Trinsoutrot-Gattin, I., Riah-Anglet, W. (2024). Which soil
microbial indicators should be included in routine laboratory tests to support
the transition to sustainable management of arable farming systems?
A meta-analysis. \emph{Ecological Indicators}, 167, 112706.

\bibitem{bhaduri2022}
Bhaduri, D., Sihi, D., Bhowmik, A., Verma, B.C., Munda, S., Dari, B. (2022).
A review on effective soil health bio-indicators for ecosystem restoration and
sustainability. \emph{Frontiers in Microbiology}, 13, 938481.

\bibitem{cardoso2013}
Cardoso, E.J.B.N., Vasconcellos, R.L.F., Bini, D., Miyauchi, M.Y.H.,
Santos, C.A., Alves, P.R.L., de Paula, A.M., Nakatani, A.S., Pereira, J.M.,
Nogueira, M.A. (2013). Soil health: looking for suitable indicators. What
should be considered to assess the effects of use and management on soil
health? \emph{Scientia Agricola}, 70(4), 274--289.

\bibitem{jia2025}
Jia, J., de Goede, R., Li, Y., Zhang, J., Wang, G., Zhang, J., Creamer, R.
(2025). Unlocking soil health: Are microbial functional genes effective
indicators? \emph{Soil Biology and Biochemistry}, 204, 109768.

\bibitem{trivedi2016}
Trivedi, P., Delgado-Baquerizo, M., Anderson, I.C., Singh, B.K. (2016).
Response of soil properties and microbial communities to agriculture:
Implications for primary productivity and soil health indicators.
\emph{Frontiers in Plant Science}, 7, 990.

\bibitem{bunemann2018}
Bünemann, E.K., Bongiorno, G., Bai, Z., Creamer, R., De Deyn, G., de Goede, R.,
et al. (2018). Soil quality --- a critical review. \emph{Soil Biology and
Biochemistry}, 120, 105--125.

\bibitem{bagnall2023}
Bagnall, D.K., Rieke, E.L., Morgan, C.L.S., Liptzin, D.L., Cappellazzi, S.B.,
Honeycutt, C.W. (2023). A minimum suite of soil health indicators for North
American agriculture. \emph{Soil Security}, 10, 100084.

\bibitem{pearsons2023}
Pearsons, K.A., et al. (2023). A tale of two systems: Does reducing tillage
affect soil health differently depending on the system? \emph{Soil \& Tillage
Research}, 226, 105562.

\bibitem{kuzyakov2010}
Kuzyakov, Y. (2010). Priming effects: Interactions between living and dead
organic matter. \emph{Soil Biology and Biochemistry}, 42(9), 1363--1371.

\bibitem{bhupenchandra2026}
Bhupenchandra, I., Sharma, L.D., Salam, M.D., Kumar, S., Chongtham, S.K.,
Choudhary, A.K., Devi, H.L., Sinyorita, S., Dutta, P., Bouket, A.C.,
Minkina, T., Keswani, C. (2026). Unravelling the role of soil enzymatic
activity as a vital bioindicator of soil health and fertility in agricultural
systems. \emph{Applied Soil Ecology}, 219, 106816.
\end{thebibliography}

\end{document}
